\section{Introduction}\label{sec:intro}

The architecture of an engineered system is chosen early, when little is known about how the system will fail, and it is chosen against criteria that are easy to quantify at that stage: performance, mass, cost, efficiency, manufacturability. Whether the failures that matter in service will be \emph{observable}---and therefore detectable before they become expensive---is rarely among those criteria. Condition monitoring is instead specified later, by different people, for an architecture that is already fixed, and its difficulties are treated as signal-processing problems to be solved with better algorithms rather than as consequences of the design decision that created them. The Design-for-X literature has produced operational methods for manufacturability, assembly, reliability and maintainability, and, in electronics, for testability \citep{Holt2010, Simpson1994}; for electromechanical systems, however, there is no equivalent method that lets a designer ask, at the point of choosing between architectures, \emph{``which of these alternatives will let me see the faults I care about, with which sensors, and how early?''}

This paper argues that the question can be answered at the design stage, and gives a method for doing so. The central idea is to treat diagnosability as a property of the triple (architecture, control structure, observation suite) rather than of the fault, and to evaluate it with a single scale-free statistic that can be computed from any set of trials---measured or simulated---in which a failure mode of known extent is introduced into an otherwise healthy system. We call the statistic the signal-to-drift ratio (SDR): the relative change of an observable feature between a reference interval and a faulted interval, divided by the size of the relative changes that the same feature exhibits between adjacent reference intervals when nothing is wrong. An SDR of one marks the point at which a fault becomes distinguishable from the ordinary non-stationarity of the healthy system at a chosen per-trial false-alarm rate. From the SDR follow the two quantities a designer needs: the \emph{minimum detectable extent} of a failure mode for a given alternative and observation suite, and the \emph{share} of operating conditions in which that extent is detectable. Both are requirement-writable, both can be compared across alternatives whose signals live in different physical units, and both can be traded against the cost of the observation suite.

The motivating domain is the generator of a large wind turbine. The choice between a permanent-magnet synchronous generator (PMSG), a squirrel-cage induction generator (SCIG) and a wound-field synchronous generator (WFSG) is one of the established architecture-level trades in turbine design \citep{Polinder2006, Polinder2013}, and electrical-subsystem failures are a major contributor to offshore downtime and unscheduled maintenance cost \citep{Tavner2007, Carroll2016}. Stator-winding insulation faults are the archetypal early-stage failure: they begin as a short-circuit across a few turns, produce subtle electrical signatures, and evolve into machine-destroying phase faults if undetected \citep{Gandhi2011, RieraGuasp2015}. A rich literature exists on detecting them in a given machine \citep{Thomson2001, Nandi2005, Bellini2008}; almost none of it asks how the choice of machine changes the problem.

A recent series of open benchmark datasets makes that question answerable empirically for the first time. \citet{Tominaga2025pmsg, Tominaga2025scig, Tominaga2024gen} subjected a PMSG, an SCIG and a variable-speed WFSG to the \emph{same} controlled inter-turn and inter-winding short-circuits, inserted at the same tapped positions of the stator winding, at matched operating points, on instrumented laboratory benches; the PMSG and SCIG share a single bench and an identical test protocol. The datasets were published for algorithm benchmarking and explicitly refrain from comparing the machines. We use them for the opposite purpose: as the test-stage instantiation of a design method whose object is the comparison itself.

The contributions are:
\begin{enumerate}
  \item A formal definition of \emph{alternative-dependent diagnosability}, with the SDR statistic, its calibration against a pooled no-fault null, a reliability screen for observables, and the derived design quantities (minimum detectable extent, detectable share, signature location).
  \item A six-step design procedure that uses these quantities to co-select a system architecture and its monitoring suite under a detectability requirement, and to assess whether a monitoring design will transfer across the alternatives of a product family.
  \item A demonstration on three generator topologies under identical winding faults, yielding four empirical findings with design content: the same design-level fault is not the same physical fault across architectures; the control structure decides where in the drive the fault becomes visible; the cheapest adequate observation suite is architecture-specific; and a learned detector transfers across architectures only when its features are referenced to the monitored unit's own healthy behaviour.
  \item A robustness and uncertainty analysis of the method's tuning choices (null quantile, window length) and of the recording-level variability behind every reported number, together with open code that reproduces all results from the public data.
\end{enumerate}

Section~\ref{sec:background} positions the work in the design and diagnosis literatures. Section~\ref{sec:method} defines the method. Section~\ref{sec:demo} describes the three alternatives, the fault protocol, the observation suites and the computation. Section~\ref{sec:results} reports the results, Section~\ref{sec:discussion} draws out the design implications and the limits of what one machine per alternative can support, and Section~\ref{sec:conclusion} concludes.
